\documentclass[12pt,]{article}
\usepackage{lmodern}
\usepackage{amssymb,amsmath}
\usepackage{ifxetex,ifluatex}
\usepackage{fixltx2e} % provides \textsubscript
\ifnum 0\ifxetex 1\fi\ifluatex 1\fi=0 % if pdftex
  \usepackage[T1]{fontenc}
  \usepackage[utf8]{inputenc}
\else % if luatex or xelatex
  \ifxetex
    \usepackage{mathspec}
  \else
    \usepackage{fontspec}
  \fi
  \defaultfontfeatures{Ligatures=TeX,Scale=MatchLowercase}
\fi
% use upquote if available, for straight quotes in verbatim environments
\IfFileExists{upquote.sty}{\usepackage{upquote}}{}
% use microtype if available
\IfFileExists{microtype.sty}{%
\usepackage{microtype}
\UseMicrotypeSet[protrusion]{basicmath} % disable protrusion for tt fonts
}{}
\usepackage[margin=1in]{geometry}
\usepackage{hyperref}
\hypersetup{unicode=true,
            pdftitle={Multidimensional Array Indexing and Storage},
            pdfauthor={Jan de Leeuw},
            pdfborder={0 0 0},
            breaklinks=true}
\urlstyle{same}  % don't use monospace font for urls
\usepackage{color}
\usepackage{fancyvrb}
\newcommand{\VerbBar}{|}
\newcommand{\VERB}{\Verb[commandchars=\\\{\}]}
\DefineVerbatimEnvironment{Highlighting}{Verbatim}{commandchars=\\\{\}}
% Add ',fontsize=\small' for more characters per line
\usepackage{framed}
\definecolor{shadecolor}{RGB}{248,248,248}
\newenvironment{Shaded}{\begin{snugshade}}{\end{snugshade}}
\newcommand{\KeywordTok}[1]{\textcolor[rgb]{0.13,0.29,0.53}{\textbf{#1}}}
\newcommand{\DataTypeTok}[1]{\textcolor[rgb]{0.13,0.29,0.53}{#1}}
\newcommand{\DecValTok}[1]{\textcolor[rgb]{0.00,0.00,0.81}{#1}}
\newcommand{\BaseNTok}[1]{\textcolor[rgb]{0.00,0.00,0.81}{#1}}
\newcommand{\FloatTok}[1]{\textcolor[rgb]{0.00,0.00,0.81}{#1}}
\newcommand{\ConstantTok}[1]{\textcolor[rgb]{0.00,0.00,0.00}{#1}}
\newcommand{\CharTok}[1]{\textcolor[rgb]{0.31,0.60,0.02}{#1}}
\newcommand{\SpecialCharTok}[1]{\textcolor[rgb]{0.00,0.00,0.00}{#1}}
\newcommand{\StringTok}[1]{\textcolor[rgb]{0.31,0.60,0.02}{#1}}
\newcommand{\VerbatimStringTok}[1]{\textcolor[rgb]{0.31,0.60,0.02}{#1}}
\newcommand{\SpecialStringTok}[1]{\textcolor[rgb]{0.31,0.60,0.02}{#1}}
\newcommand{\ImportTok}[1]{#1}
\newcommand{\CommentTok}[1]{\textcolor[rgb]{0.56,0.35,0.01}{\textit{#1}}}
\newcommand{\DocumentationTok}[1]{\textcolor[rgb]{0.56,0.35,0.01}{\textbf{\textit{#1}}}}
\newcommand{\AnnotationTok}[1]{\textcolor[rgb]{0.56,0.35,0.01}{\textbf{\textit{#1}}}}
\newcommand{\CommentVarTok}[1]{\textcolor[rgb]{0.56,0.35,0.01}{\textbf{\textit{#1}}}}
\newcommand{\OtherTok}[1]{\textcolor[rgb]{0.56,0.35,0.01}{#1}}
\newcommand{\FunctionTok}[1]{\textcolor[rgb]{0.00,0.00,0.00}{#1}}
\newcommand{\VariableTok}[1]{\textcolor[rgb]{0.00,0.00,0.00}{#1}}
\newcommand{\ControlFlowTok}[1]{\textcolor[rgb]{0.13,0.29,0.53}{\textbf{#1}}}
\newcommand{\OperatorTok}[1]{\textcolor[rgb]{0.81,0.36,0.00}{\textbf{#1}}}
\newcommand{\BuiltInTok}[1]{#1}
\newcommand{\ExtensionTok}[1]{#1}
\newcommand{\PreprocessorTok}[1]{\textcolor[rgb]{0.56,0.35,0.01}{\textit{#1}}}
\newcommand{\AttributeTok}[1]{\textcolor[rgb]{0.77,0.63,0.00}{#1}}
\newcommand{\RegionMarkerTok}[1]{#1}
\newcommand{\InformationTok}[1]{\textcolor[rgb]{0.56,0.35,0.01}{\textbf{\textit{#1}}}}
\newcommand{\WarningTok}[1]{\textcolor[rgb]{0.56,0.35,0.01}{\textbf{\textit{#1}}}}
\newcommand{\AlertTok}[1]{\textcolor[rgb]{0.94,0.16,0.16}{#1}}
\newcommand{\ErrorTok}[1]{\textcolor[rgb]{0.64,0.00,0.00}{\textbf{#1}}}
\newcommand{\NormalTok}[1]{#1}
\usepackage{graphicx,grffile}
\makeatletter
\def\maxwidth{\ifdim\Gin@nat@width>\linewidth\linewidth\else\Gin@nat@width\fi}
\def\maxheight{\ifdim\Gin@nat@height>\textheight\textheight\else\Gin@nat@height\fi}
\makeatother
% Scale images if necessary, so that they will not overflow the page
% margins by default, and it is still possible to overwrite the defaults
% using explicit options in \includegraphics[width, height, ...]{}
\setkeys{Gin}{width=\maxwidth,height=\maxheight,keepaspectratio}
\IfFileExists{parskip.sty}{%
\usepackage{parskip}
}{% else
\setlength{\parindent}{0pt}
\setlength{\parskip}{6pt plus 2pt minus 1pt}
}
\setlength{\emergencystretch}{3em}  % prevent overfull lines
\providecommand{\tightlist}{%
  \setlength{\itemsep}{0pt}\setlength{\parskip}{0pt}}
\setcounter{secnumdepth}{5}
% Redefines (sub)paragraphs to behave more like sections
\ifx\paragraph\undefined\else
\let\oldparagraph\paragraph
\renewcommand{\paragraph}[1]{\oldparagraph{#1}\mbox{}}
\fi
\ifx\subparagraph\undefined\else
\let\oldsubparagraph\subparagraph
\renewcommand{\subparagraph}[1]{\oldsubparagraph{#1}\mbox{}}
\fi

%%% Use protect on footnotes to avoid problems with footnotes in titles
\let\rmarkdownfootnote\footnote%
\def\footnote{\protect\rmarkdownfootnote}

%%% Change title format to be more compact
\usepackage{titling}

% Create subtitle command for use in maketitle
\providecommand{\subtitle}[1]{
  \posttitle{
    \begin{center}\large#1\end{center}
    }
}

\setlength{\droptitle}{-2em}

  \title{Multidimensional Array Indexing and Storage}
    \pretitle{\vspace{\droptitle}\centering\huge}
  \posttitle{\par}
    \author{Jan de Leeuw}
    \preauthor{\centering\large\emph}
  \postauthor{\par}
      \predate{\centering\large\emph}
  \postdate{\par}
    \date{Version 27, August 31, 2017}


\begin{document}
\maketitle
\begin{abstract}
We give R, C, and R-\textgreater{}C code to access lineary stored
multidimensional arrays and compactly stored multidimensional
super-symmetric arrays.
\end{abstract}

{
\setcounter{tocdepth}{3}
\tableofcontents
}
\begin{center}\rule{0.5\linewidth}{\linethickness}\end{center}

Note: This is a working paper which will be expanded/updated frequently.
All suggestions for improvement are welcome. The directory
\href{http://deleeuwpdx.net/pubfolders/indexing}{deleeuwpdx.net/pubfolders/indexing}
has a pdf version, the bib file, the complete Rmd file with the code
chunks, and the R and C source code.

\section{Introduction}\label{introduction}

Suppose \(A\) is an array of \emph{rank} \(r\) and \emph{dimension}
\((n_1,\cdots,n_r)\). Here rank is used in the APL sense, i.e.~the
length of the dimension vector.

\(A\) has elements \(a_{i_1\cdots i_m}\). It is easy enough to
visualize, or at least conceptualize, these \(n:=n_1\times\cdots n_m\)
elements in an \(m\)-dimensional rectangular block of cells, but
unfortunately such blocks cannot be stored directly in contiguous
locations in computer memory. The array \(A\) is usually stored in \(n\)
contiguous locations in memory, i.e.~in a form that corresponds with a
vector, an array of rank one.

Say \(b\) is a pointer to the first element stored. We need a mapping
\(f\) from \(\mathcal{I}_{n_1}\otimes\cdots\mathcal{I}_{n_m}\) into
\(\mathcal{I}_n\) such that \(b+f(i_1,\cdots,i_m)\) is a pointer to
\(a_{i_1\cdots i_m}\). Or, using C,
\(b[f(i_1,\cdots,i_m)]=*(b+f(i_1,\cdots,i_m))\) must be equal to
\(a_{i_1\cdots i_m}\). Let's call such an \(f\) an \emph{index mapping}.
We are interested in both the index mapping \(f\) and its inverse
\(f^{-1}:\mathcal{I}_n\rightarrow\mathcal{I}_{n_1}\otimes\cdots\mathcal{I}_{n_m}\).

If \(A\) is symmetric, or anti-symmetric, or triangular, or hollow, then
it is not necessary to store all its elements. This means that \(f\) is
no longer bijective, because fewer than \(n\) elements are stored in
memory. In fact, the main motivation for this paper is to derive the
index mapping, and its inverse, for the super-symmetric arrays of
partials derivatives discussed in De Leeuw (2017).

In this paper I review and program index mappings for arrays that can be
used to assess and retrieve array elements from contiguous memory
locations. I am \(99.99\%\) sure there is nothing new in the paper, I
just wanted to collect these results in a single convenient place, and
provide code. As for the programming languages, we first develop in R,
and then translate to C.

\section{Matrices}\label{matrices}

\subsection{General Rectangular
Matrices}\label{general-rectangular-matrices}

Storing the elements of a general rectangular matrix in contiguous
locations in memory can be done in many ways, but two specific ones are
the most obvious. We can store row one first, then continue with row
two, and so on. This is \emph{row major order}. Or we can store column
one, append column two, and so on. That is \emph{column major order}. If
\(A\) is

\begin{verbatim}
##      [,1] [,2] [,3]
## [1,]   3    2    1 
## [2,]  10    6    7 
## [3,]   8   12    5 
## [4,]  11    9    4
\end{verbatim}

then column major order stores the matrix as

\begin{verbatim}
##  [1]   3  10   8  11   2   6  12   9   1   7   5   4
\end{verbatim}

and row major order stores it as

\begin{verbatim}
##  [1]   3   2   1  10   6   7   8  12   5  11   9   4
\end{verbatim}

The index mapping for column major order is

\begin{verbatim}
##  1  1  ***   1 
##  2  1  ***   2 
##  3  1  ***   3 
##  4  1  ***   4 
##  1  2  ***   5 
##  2  2  ***   6 
##  3  2  ***   7 
##  4  2  ***   8 
##  1  3  ***   9 
##  2  3  ***  10 
##  3  3  ***  11 
##  4  3  ***  12
\end{verbatim}

and for row major order it is

\begin{verbatim}
##  1  1  ***   1 
##  1  2  ***   2 
##  1  3  ***   3 
##  2  1  ***   4 
##  2  2  ***   5 
##  2  3  ***   6 
##  3  1  ***   7 
##  3  2  ***   8 
##  3  3  ***   9 
##  4  1  ***  10 
##  4  2  ***  11 
##  4  3  ***  12
\end{verbatim}

Because the row index changes fastest, we could call column major
ordering \emph{first-fast}. Row major ordering has the second, i.e.~the
last index changing fastest, and could therefore be called
\emph{last-fast}. That terminology has the advantage it generalizes to
multidimensional arrays, while the use of rows and columns is more or
less limited to matrices.

The main application we have in mind is to pass matrices and arrays from
R to C, with I/O and memory allocation done in R, and computation in C.
This implies that the C versions of our indexing routines are the most
useful ones.

R stores matrices first-fast (showing its origins, as a REPL interface
to FORTRAN routines). Thus then \texttt{a{[}3{]}\ =}8 and
\texttt{a{[}11{]}\ =}5, because \texttt{as.vector(a)} is

\begin{verbatim}
##  [1]  3 10  8 11  2  6 12  9  1  7  5  4
\end{verbatim}

First-fast is how the arrays arrive in C, where the index calculations
are done to retrieve the elements. As a consequence, both for matrices
and arrays, first-fast is the most relevant storage mode for our
purposes. We do not have to worry about how C stores matrices
multidimensional arrays, because for our purposes it is better to
pretend that C does not have any matrices or multidimensional arrays at
all.

\subsection{Symmetric and Triangular
Matrices}\label{symmetric-and-triangular-matrices}

For symmetric and anti-symmetric matrices we only have to store half of
the elements. This means making two choices, first if we want to store
the upper or the lower triangle, second if we want to use first-fast or
last-fast. Storing the lower triangle means storing all \(a_{ij}\) with
\(i\geq j\), storing the upper triangle stores all \(a_{ij}\) with
\(i\leq j\). Because of generalization to multidimensional arrays we
call the first option \emph{decreasing} and the second option
\emph{increasing}.

We will choose the upper or increasing triangle combined with first-fast
storage throughout. For a symmetric matrix we only have to store
\(N=\frac12 n(n+1)\) elements. The mapping function will be defined for
all \(i\leq j\), and it maps those index pairs bijectively on
\(1,\cdots,N\). This mapping function also has an inverse. We then
extend the mapping function to all index pairs by first ordering the
pair, making the extended mapping function surjective. We will not give
the mapping function yet, because it will be just a special case of the
one for super-symmetric arrays.

Just as an aside, our conventions imply that the elements of a lower
triangular matrix, in which the non-zero \(a_{ij}\) have \(i\geq j\),
will be stored in increasing or upper triangle form, in other words as
its transpose. The extended mapping function for that case, which we do
not really consider in detail, will have to take that into account.
Also, strictly lower and upper triangular matrices, anti-symmetric
matrices, and symmetric hollow matrices such as distance matrices, only
need to store \(\frac12 n(n-1)\) elements. We do not discuss storage
schemes that leave out the diagonal. We always store the diagonal
elements, but the extended mapping function has no need to access them.

\section{Arrays}\label{arrays}

\subsection{General Arrays}\label{general-arrays}

Linear storage of the elements of a general array or rank \(m\) and
dimension \((n_1,\cdots,n_m)\) can be done in many ways. First-fast and
last-fast are just two special cases of the \(m!\) ways in which we can
nest the index loops. Since going through all possibilities is both
boring and unnessary we limit ourselves to first-fast. If, for whatever
reason, another permutation of the indices is needed we can first
transpose the array in R, for instance using \texttt{aperm()} from the
base package, or \texttt{aplTranspose()} from De Leeuw and Yajima
(2016).

The R version of the mapping function is \texttt{fArrayFirst()}. Note
that it is identical to the \texttt{aplDecode()} function from De Leeuw
and Yajima (2016). For completeness we also give
\texttt{fArrayFirstInverse()}, which is \texttt{aplEncode()}, the
inverse \(f^{-1}\), mapping \(\mathcal{I}_n\) into
\(\mathcal{I}_{n_1}\otimes\cdots\mathcal{I}_{n_m}\). Our example is a
\(4\times 3\times 2\) array. In tabular form the index mapping is

\begin{verbatim}
##  1  1  1  ***   1 
##  2  1  1  ***   2 
##  3  1  1  ***   3 
##  4  1  1  ***   4 
##  1  2  1  ***   5 
##  2  2  1  ***   6 
##  3  2  1  ***   7 
##  4  2  1  ***   8 
##  1  3  1  ***   9 
##  2  3  1  ***  10 
##  3  3  1  ***  11 
##  4  3  1  ***  12 
##  1  1  2  ***  13 
##  2  1  2  ***  14 
##  3  1  2  ***  15 
##  4  1  2  ***  16 
##  1  2  2  ***  17 
##  2  2  2  ***  18 
##  3  2  2  ***  19 
##  4  2  2  ***  20 
##  1  3  2  ***  21 
##  2  3  2  ***  22 
##  3  3  2  ***  23 
##  4  3  2  ***  24
\end{verbatim}

and in array form it is

\begin{verbatim}
## , , 1
## 
##      [,1] [,2] [,3]
## [1,]    1    5    9
## [2,]    2    6   10
## [3,]    3    7   11
## [4,]    4    8   12
## 
## , , 2
## 
##      [,1] [,2] [,3]
## [1,]   13   17   21
## [2,]   14   18   22
## [3,]   15   19   23
## [4,]   16   20   24
\end{verbatim}

\subsection{Super-symmetry}\label{super-symmetry}

An array \(A\) of rank \(m\) and dimension \((n,\cdots,n)\) is
\emph{super-symmetric} if \(a_{i_1\cdots i_m}\) is invariant under
permutation of indices. Think of arrays of product moments, cumulants,
or partial derivatives. Super-symmetry generalizes symmetry of matrices.

In order not to drown in the \((m!)\times (m!)\) combinations of
what-fast and which-triangle, we limit ourselves to first-fast and
increasing (i.e.~index tuples with \(i_1\leq\cdots\leq i_m\)).

The number of elements we have to store is \(\binom{m+n-1}{m}\), instead
of \(n^m\). This can be a considerable saving of storage if \(m\) and/or
\(n\) are at all large. For \(n=10\) and \(m=4\), for example, we save
\(93\%\) of the storage. For large \(n\) the savings are approximately
\(100*(1-\frac{1}{m!})\%\).

The mapping function for a super-symmetric array, with first-fast and
increasing index vectors, is \texttt{fSupSymIncreasingFirst()}. Note
that this function does not need to know the dimensions of the array.
Also it starts by sorting the elements of the index vector. Thus

\begin{Shaded}
\begin{Highlighting}[]
\KeywordTok{fSupSymIncreasingFirstRC}\NormalTok{(}\KeywordTok{c}\NormalTok{(}\DecValTok{1}\NormalTok{,}\DecValTok{2}\NormalTok{,}\DecValTok{2}\NormalTok{,}\DecValTok{3}\NormalTok{))}
\end{Highlighting}
\end{Shaded}

\begin{verbatim}
## [1] 8
\end{verbatim}

and also

\begin{Shaded}
\begin{Highlighting}[]
\KeywordTok{fSupSymIncreasingFirstRC}\NormalTok{(}\KeywordTok{c}\NormalTok{(}\DecValTok{2}\NormalTok{,}\DecValTok{1}\NormalTok{,}\DecValTok{3}\NormalTok{,}\DecValTok{2}\NormalTok{))}
\end{Highlighting}
\end{Shaded}

\begin{verbatim}
## [1] 8
\end{verbatim}

The inverse mapping is \texttt{fSupSymIncreasingFirstInverse()}.

As an example, here is the result for a four-dimensional super-symmetric
array of order four. There are only 35, instead of 256, elements to
consider.

\begin{verbatim}
##  1  1  1  1  ***   1 
##  1  1  1  2  ***   2 
##  1  1  2  2  ***   3 
##  1  2  2  2  ***   4 
##  2  2  2  2  ***   5 
##  1  1  1  3  ***   6 
##  1  1  2  3  ***   7 
##  1  2  2  3  ***   8 
##  2  2  2  3  ***   9 
##  1  1  3  3  ***  10 
##  1  2  3  3  ***  11 
##  2  2  3  3  ***  12 
##  1  3  3  3  ***  13 
##  2  3  3  3  ***  14 
##  3  3  3  3  ***  15 
##  1  1  1  4  ***  16 
##  1  1  2  4  ***  17 
##  1  2  2  4  ***  18 
##  2  2  2  4  ***  19 
##  1  1  3  4  ***  20 
##  1  2  3  4  ***  21 
##  2  2  3  4  ***  22 
##  1  3  3  4  ***  23 
##  2  3  3  4  ***  24 
##  3  3  3  4  ***  25 
##  1  1  4  4  ***  26 
##  1  2  4  4  ***  27 
##  2  2  4  4  ***  28 
##  1  3  4  4  ***  29 
##  2  3  4  4  ***  30 
##  3  3  4  4  ***  31 
##  1  4  4  4  ***  32 
##  2  4  4  4  ***  33 
##  3  4  4  4  ***  34 
##  4  4  4  4  ***  35
\end{verbatim}

\section{Implementation}\label{implementation}

There are four versions of the two basic functions
\texttt{fArrayFirst()} for general arrays and
\texttt{fSupSymIncreasingFirst()} for super-symmetric arrays. Two are in
C, one which passes by value and returns integers, and a second one
which wraps the first one, and passes by reference and returns void. And
two are in R, one in pure R, and one which uses .C() to wrap a call to
the compiled C code.

The two inverse functions \texttt{fArrayFirstInverse()} and
\texttt{fSupSymIncreasingFirstInverse()} have only a single C version,
which passes by reference and returns void. There are again two R
versions of each, one in pure R and one a wrapper around a .C() call to
the compiled C code.

\section{Appendix: Code}\label{appendix-code}

\subsection{R Code}\label{r-code}

\subsubsection{indexing.R}\label{indexing.r}

\begin{Shaded}
\begin{Highlighting}[]
\NormalTok{fArrayFirstR <-}\StringTok{ }\ControlFlowTok{function}\NormalTok{ (cell, dimension) \{}
\NormalTok{  rank <-}\StringTok{ }\KeywordTok{length}\NormalTok{ (cell)}
\NormalTok{  k <-}\StringTok{ }\KeywordTok{cumprod}\NormalTok{ (}\KeywordTok{c}\NormalTok{(}\DecValTok{1}\NormalTok{, dimension))[}\DecValTok{1}\OperatorTok{:}\NormalTok{rank]}
  \KeywordTok{return}\NormalTok{ (}\DecValTok{1} \OperatorTok{+}\StringTok{ }\KeywordTok{sum}\NormalTok{ ((cell }\OperatorTok{-}\StringTok{ }\DecValTok{1}\NormalTok{) }\OperatorTok{*}\StringTok{ }\NormalTok{k))}
\NormalTok{\}}

\NormalTok{fSupSymIncreasingFirstR <-}\StringTok{ }\ControlFlowTok{function}\NormalTok{ (cell) \{}
\NormalTok{  f <-}\StringTok{ }\ControlFlowTok{function}\NormalTok{ (r)}
    \KeywordTok{choose}\NormalTok{ (r }\OperatorTok{+}\StringTok{ }\NormalTok{(cell[r] }\OperatorTok{-}\StringTok{ }\DecValTok{1}\NormalTok{) }\OperatorTok{-}\StringTok{ }\DecValTok{1}\NormalTok{, r)}
\NormalTok{  rank <-}\StringTok{ }\KeywordTok{length}\NormalTok{ (cell)}
\NormalTok{  cell <-}\StringTok{ }\KeywordTok{sort}\NormalTok{ (cell)}
  \KeywordTok{return}\NormalTok{ (}\DecValTok{1} \OperatorTok{+}\StringTok{ }\KeywordTok{sum}\NormalTok{ (}\KeywordTok{sapply}\NormalTok{ (}\DecValTok{1}\OperatorTok{:}\NormalTok{rank, f)))}
\NormalTok{\}}

\NormalTok{fArrayFirstInverseR <-}\StringTok{ }\ControlFlowTok{function}\NormalTok{(index, dimension) \{}
\NormalTok{  rank <-}\StringTok{ }\KeywordTok{length}\NormalTok{ (dimension)}
\NormalTok{  b <-}\StringTok{ }\KeywordTok{cumprod}\NormalTok{ (}\KeywordTok{c}\NormalTok{(}\DecValTok{1}\NormalTok{, dimension))[}\DecValTok{1}\OperatorTok{:}\NormalTok{rank]}
\NormalTok{  r <-}\StringTok{ }\KeywordTok{rep}\NormalTok{(}\DecValTok{0}\NormalTok{, }\KeywordTok{length}\NormalTok{(b))}
\NormalTok{  s <-}\StringTok{ }\NormalTok{index }\OperatorTok{-}\StringTok{ }\DecValTok{1}
  \ControlFlowTok{for}\NormalTok{ (j }\ControlFlowTok{in}\NormalTok{ rank}\OperatorTok{:}\DecValTok{1}\NormalTok{) \{}
\NormalTok{    r[j] <-}\StringTok{ }\NormalTok{s }\OperatorTok{%/%}\StringTok{ }\NormalTok{b[j]}
\NormalTok{    s <-}\StringTok{ }\NormalTok{s }\OperatorTok{-}\StringTok{ }\NormalTok{r[j] }\OperatorTok{*}\StringTok{ }\NormalTok{b[j]}
\NormalTok{  \}}
  \KeywordTok{return}\NormalTok{(}\DecValTok{1} \OperatorTok{+}\StringTok{ }\NormalTok{r)}
\NormalTok{\}}

\NormalTok{fSupSymIncreasingFirstInverseR <-}\StringTok{ }\ControlFlowTok{function}\NormalTok{ (dimension, rank, index) \{}
\NormalTok{  last.true <-}\StringTok{ }\ControlFlowTok{function}\NormalTok{ (x) \{}
\NormalTok{    w <-}\StringTok{ }\KeywordTok{which}\NormalTok{ (x)}
    \KeywordTok{return}\NormalTok{ (w[}\KeywordTok{length}\NormalTok{(w)])}
\NormalTok{  \}}
\NormalTok{  a <-}\StringTok{ }\KeywordTok{rep}\NormalTok{ (}\DecValTok{0}\NormalTok{, rank)}
\NormalTok{  v <-}\StringTok{ }\NormalTok{index }\OperatorTok{-}\StringTok{ }\DecValTok{1}
  \ControlFlowTok{for}\NormalTok{ (k }\ControlFlowTok{in}\NormalTok{ rank}\OperatorTok{:}\DecValTok{1}\NormalTok{) \{}
\NormalTok{    s <-}\StringTok{ }\KeywordTok{choose}\NormalTok{ (k }\OperatorTok{+}\StringTok{ }\NormalTok{(}\DecValTok{0}\OperatorTok{:}\NormalTok{(dimension }\OperatorTok{-}\StringTok{ }\DecValTok{1}\NormalTok{)) }\OperatorTok{-}\StringTok{ }\DecValTok{1}\NormalTok{, k)}
\NormalTok{    u <-}\StringTok{ }\KeywordTok{last.true}\NormalTok{ (v }\OperatorTok{>=}\StringTok{ }\NormalTok{s)}
\NormalTok{    a[k] <-}\StringTok{ }\NormalTok{u}
\NormalTok{    v <-}\StringTok{ }\NormalTok{v }\OperatorTok{-}\StringTok{ }\NormalTok{s[u]}
\NormalTok{  \}}
  \KeywordTok{return}\NormalTok{ (a)}
\NormalTok{\}}
\end{Highlighting}
\end{Shaded}

\subsubsection{indexingC.R}\label{indexingc.r}

\begin{Shaded}
\begin{Highlighting}[]
\KeywordTok{dyn.load}\NormalTok{(}\StringTok{"indexing.so"}\NormalTok{)}

\NormalTok{fArrayFirstRC <-}\StringTok{ }\ControlFlowTok{function}\NormalTok{ (cell, dimension) \{}
\NormalTok{  rank <-}\StringTok{ }\KeywordTok{length}\NormalTok{ (cell)}
\NormalTok{  h <-}
\StringTok{    }\KeywordTok{.C}\NormalTok{(}
      \StringTok{"fArrayFirstGlue"}\NormalTok{,}
      \KeywordTok{as.integer}\NormalTok{(cell),}
      \KeywordTok{as.integer}\NormalTok{(dimension),}
      \KeywordTok{as.integer}\NormalTok{(rank),}
      \DataTypeTok{index =} \KeywordTok{as.integer}\NormalTok{(}\DecValTok{0}\NormalTok{)}
\NormalTok{    )}
  \KeywordTok{return}\NormalTok{ (h}\OperatorTok{$}\NormalTok{index)}
\NormalTok{\}}


\NormalTok{fSupSymIncreasingFirstRC <-}\StringTok{ }\ControlFlowTok{function}\NormalTok{ (cell) \{}
\NormalTok{  rank <-}\StringTok{ }\KeywordTok{length}\NormalTok{(cell)}
\NormalTok{  h <-}
\StringTok{    }\KeywordTok{.C}\NormalTok{(}\StringTok{"fSupSymIncreasingGlue"}\NormalTok{,}
       \KeywordTok{as.integer}\NormalTok{(cell),}
       \KeywordTok{as.integer}\NormalTok{(rank),}
       \DataTypeTok{index =} \KeywordTok{as.integer}\NormalTok{(}\DecValTok{0}\NormalTok{))}
  \KeywordTok{return}\NormalTok{ (h}\OperatorTok{$}\NormalTok{index)}
\NormalTok{\}}

\NormalTok{fArrayFirstInverseRC <-}\StringTok{ }\ControlFlowTok{function}\NormalTok{(index, dimension) \{}
\NormalTok{  rank <-}\StringTok{ }\KeywordTok{length}\NormalTok{ (dimension)}
\NormalTok{  h <-}
\StringTok{    }\KeywordTok{.C}\NormalTok{(}
      \StringTok{"fArrayFirstInverse"}\NormalTok{,}
      \KeywordTok{as.integer}\NormalTok{(rank),}
      \KeywordTok{as.integer}\NormalTok{ (dimension),}
      \KeywordTok{as.integer}\NormalTok{(index),}
      \DataTypeTok{cell =} \KeywordTok{as.integer}\NormalTok{(}\KeywordTok{rep}\NormalTok{(}\DecValTok{0}\NormalTok{, rank))}
\NormalTok{    )}
  \KeywordTok{return}\NormalTok{ (h}\OperatorTok{$}\NormalTok{cell)}
\NormalTok{\}}

\NormalTok{fSupSymIncreasingFirstInverseRC <-}\StringTok{ }\ControlFlowTok{function}\NormalTok{ (dimension, rank, index) \{}
\NormalTok{  h <-}
\StringTok{    }\KeywordTok{.C}\NormalTok{(}
      \StringTok{"fSupSymIncreasingFirstInverse"}\NormalTok{,}
      \KeywordTok{as.integer}\NormalTok{(dimension),}
      \KeywordTok{as.integer}\NormalTok{(rank),}
      \KeywordTok{as.integer}\NormalTok{(index),}
      \DataTypeTok{cell =} \KeywordTok{as.integer}\NormalTok{(}\KeywordTok{rep}\NormalTok{(}\DecValTok{0}\NormalTok{, rank))}
\NormalTok{    )}
  \KeywordTok{return}\NormalTok{ (h}\OperatorTok{$}\NormalTok{cell)}
\NormalTok{\}}
\end{Highlighting}
\end{Shaded}

\subsection{C Code}\label{c-code}

\subsubsection{indexing.h}\label{indexing.h}

\begin{Shaded}
\begin{Highlighting}[]
\PreprocessorTok{#ifndef INDEXING_H}
\PreprocessorTok{#define INDEXING_H}

\PreprocessorTok{#include }\ImportTok{<stdlib.h>}

\DataTypeTok{static} \KeywordTok{inline} \DataTypeTok{int}\NormalTok{ IMIN(}\DataTypeTok{const} \DataTypeTok{int}\NormalTok{ x, }\DataTypeTok{const} \DataTypeTok{int}\NormalTok{ y) \{ }\ControlFlowTok{return}\NormalTok{ (x < y) ? x : y; \}}
\DataTypeTok{static} \KeywordTok{inline} \DataTypeTok{int}\NormalTok{ VINDEX(}\DataTypeTok{const} \DataTypeTok{int}\NormalTok{ i) \{ }\ControlFlowTok{return}\NormalTok{ i - }\DecValTok{1}\NormalTok{; \}}

\PreprocessorTok{#endif }\CommentTok{/* INDEXING_H */}
\end{Highlighting}
\end{Shaded}

\subsubsection{indexing.c}\label{indexing.c}

\begin{Shaded}
\begin{Highlighting}[]

\PreprocessorTok{#include }\ImportTok{"indexing.h"}

\DataTypeTok{int}\NormalTok{ binCoef(}\DataTypeTok{const} \DataTypeTok{int}\NormalTok{ n, }\DataTypeTok{const} \DataTypeTok{int}\NormalTok{ m) \{}
    \DataTypeTok{int}\NormalTok{ *work = (}\DataTypeTok{int}\NormalTok{ *)calloc((}\DataTypeTok{size_t}\NormalTok{)m + }\DecValTok{1}\NormalTok{, }\KeywordTok{sizeof}\NormalTok{(}\DataTypeTok{int}\NormalTok{));}
\NormalTok{    work[}\DecValTok{0}\NormalTok{] = }\DecValTok{1}\NormalTok{;}
    \ControlFlowTok{for}\NormalTok{ (}\DataTypeTok{int}\NormalTok{ i = }\DecValTok{1}\NormalTok{; i <= n; i++) \{}
        \ControlFlowTok{for}\NormalTok{ (}\DataTypeTok{int}\NormalTok{ j = IMIN(i, m); j > }\DecValTok{0}\NormalTok{; j--) \{}
\NormalTok{            work[j] = work[j] + work[j - }\DecValTok{1}\NormalTok{];}
\NormalTok{        \}}
\NormalTok{    \}}
    \DataTypeTok{int}\NormalTok{ choose = work[m];}
\NormalTok{    free(work);}
    \ControlFlowTok{return}\NormalTok{ (choose);}
\NormalTok{\}}

\DataTypeTok{int}\NormalTok{ int_cmp(}\DataTypeTok{const} \DataTypeTok{void}\NormalTok{ *x, }\DataTypeTok{const} \DataTypeTok{void}\NormalTok{ *y) \{}
    \DataTypeTok{const} \DataTypeTok{int}\NormalTok{ *ix = (}\DataTypeTok{const} \DataTypeTok{int}\NormalTok{ *)x;}
    \DataTypeTok{const} \DataTypeTok{int}\NormalTok{ *iy = (}\DataTypeTok{const} \DataTypeTok{int}\NormalTok{ *)y;}
    \ControlFlowTok{return}\NormalTok{ (*ix - *iy);}
\NormalTok{\}}

\DataTypeTok{int}\NormalTok{ fArrayFirst(}\DataTypeTok{const} \DataTypeTok{int}\NormalTok{ *cell, }\DataTypeTok{const} \DataTypeTok{int}\NormalTok{ *dimension, }\DataTypeTok{const} \DataTypeTok{int}\NormalTok{ rank) \{}
    \DataTypeTok{int}\NormalTok{ aux = }\DecValTok{1}\NormalTok{, result = }\DecValTok{1}\NormalTok{;}
    \ControlFlowTok{for}\NormalTok{ (}\DataTypeTok{int}\NormalTok{ i = }\DecValTok{1}\NormalTok{; i <= rank; i++) \{}
\NormalTok{        result += (cell[VINDEX(i)] - }\DecValTok{1}\NormalTok{) * aux;}
\NormalTok{        aux *= dimension[VINDEX(i)];}
\NormalTok{    \}}
    \ControlFlowTok{return}\NormalTok{ (result);}
\NormalTok{\}}

\DataTypeTok{int}\NormalTok{ fSupSymIncreasing(}\DataTypeTok{int}\NormalTok{ *cell, }\DataTypeTok{const} \DataTypeTok{int}\NormalTok{ rank) \{}
    \DataTypeTok{int}\NormalTok{ result = }\DecValTok{1}\NormalTok{;}
\NormalTok{    (}\DataTypeTok{void}\NormalTok{)qsort(cell, (}\DataTypeTok{size_t}\NormalTok{)rank, }\KeywordTok{sizeof}\NormalTok{(}\DataTypeTok{int}\NormalTok{), int_cmp);}
    \ControlFlowTok{for}\NormalTok{ (}\DataTypeTok{int}\NormalTok{ i = }\DecValTok{1}\NormalTok{; i <= rank; i++) \{}
        \DataTypeTok{int}\NormalTok{ k = i + (cell[VINDEX(i)] - }\DecValTok{1}\NormalTok{) - }\DecValTok{1}\NormalTok{;}
\NormalTok{        result += binCoef(k, i);}
\NormalTok{    \}}
    \ControlFlowTok{return}\NormalTok{ (result);}
\NormalTok{\}}

\DataTypeTok{void}\NormalTok{ fArrayFirstInverse(}\DataTypeTok{const} \DataTypeTok{int}\NormalTok{ *rank, }\DataTypeTok{const} \DataTypeTok{int}\NormalTok{ *dimension, }\DataTypeTok{const} \DataTypeTok{int}\NormalTok{ *index,}
                        \DataTypeTok{int}\NormalTok{ *cell) \{}
    \DataTypeTok{int}\NormalTok{ m = *rank, l = *index, b = }\DecValTok{1}\NormalTok{;}
    \ControlFlowTok{for}\NormalTok{ (}\DataTypeTok{int}\NormalTok{ j = }\DecValTok{1}\NormalTok{; j < m; j++) \{}
\NormalTok{        b *= dimension[VINDEX(j)];}
\NormalTok{    \}}
    \ControlFlowTok{for}\NormalTok{ (}\DataTypeTok{int}\NormalTok{ j = m - }\DecValTok{1}\NormalTok{; j > }\DecValTok{0}\NormalTok{; j--) \{}
        \DataTypeTok{int}\NormalTok{ k = (l - }\DecValTok{1}\NormalTok{) / b;}
\NormalTok{        cell[j] = k + }\DecValTok{1}\NormalTok{;}
\NormalTok{        l -= k * b;}
\NormalTok{        b /= dimension[j - }\DecValTok{1}\NormalTok{];}
\NormalTok{    \}}
\NormalTok{    cell[VINDEX(}\DecValTok{1}\NormalTok{)] = l;}
    \ControlFlowTok{return}\NormalTok{;}
\NormalTok{\}}

\DataTypeTok{void}\NormalTok{ fSupSymIncreasingFirstInverse(}\DataTypeTok{const} \DataTypeTok{int}\NormalTok{ *dimension, }\DataTypeTok{const} \DataTypeTok{int}\NormalTok{ *rank,}
                                   \DataTypeTok{const} \DataTypeTok{int}\NormalTok{ *index, }\DataTypeTok{int}\NormalTok{ *cell) \{}
    \DataTypeTok{int}\NormalTok{ n = *dimension, m = *rank, l = *index, v = l - }\DecValTok{1}\NormalTok{;}
    \ControlFlowTok{for}\NormalTok{ (}\DataTypeTok{int}\NormalTok{ k = m; k >= }\DecValTok{1}\NormalTok{; k--) \{}
        \ControlFlowTok{for}\NormalTok{ (}\DataTypeTok{int}\NormalTok{ j = }\DecValTok{0}\NormalTok{; j < n; j++) \{}
            \DataTypeTok{int}\NormalTok{ sj = binCoef(k + j - }\DecValTok{1}\NormalTok{, k), sk = binCoef(k + j, k);}
            \ControlFlowTok{if}\NormalTok{ (v < sk) \{}
\NormalTok{                cell[VINDEX(k)] = j + }\DecValTok{1}\NormalTok{;}
\NormalTok{                v -= sj;}
                \ControlFlowTok{break}\NormalTok{;}
\NormalTok{            \}}
\NormalTok{        \}}
\NormalTok{    \}}
    \ControlFlowTok{return}\NormalTok{;}
\NormalTok{\}}

\DataTypeTok{void}\NormalTok{ binCoefGlue(}\DataTypeTok{const} \DataTypeTok{int}\NormalTok{ *n, }\DataTypeTok{const} \DataTypeTok{int}\NormalTok{ *m, }\DataTypeTok{int}\NormalTok{ *choose) \{}
\NormalTok{    *choose = binCoef(*n, *m);}
\NormalTok{\}}

\DataTypeTok{void}\NormalTok{ fArrayFirstGlue(}\DataTypeTok{const} \DataTypeTok{int}\NormalTok{ *cell, }\DataTypeTok{const} \DataTypeTok{int}\NormalTok{ *dimension, }\DataTypeTok{const} \DataTypeTok{int}\NormalTok{ *rank,}
                     \DataTypeTok{int}\NormalTok{ *result) \{}
\NormalTok{    *result = fArrayFirst(cell, dimension, *rank);}
\NormalTok{\}}

\DataTypeTok{void}\NormalTok{ fSupSymIncreasingGlue(}\DataTypeTok{int}\NormalTok{ *cell, }\DataTypeTok{const} \DataTypeTok{int}\NormalTok{ *rank, }\DataTypeTok{int}\NormalTok{ *result) \{}
\NormalTok{    *result = fSupSymIncreasing(cell, *rank);}
\NormalTok{\}}
\end{Highlighting}
\end{Shaded}

\section*{References}\label{references}
\addcontentsline{toc}{section}{References}

\hypertarget{refs}{}
\hypertarget{ref-deleeuw_E_17r}{}
De Leeuw, J. 2017. ``Higher Partials of fStress. Who Needs Them ?''
\url{http://deleeuwpdx.net/pubfolders/fStress/fStress.pdf}.

\hypertarget{ref-deleeuw_yajima_E_16}{}
De Leeuw, J., and M. Yajima. 2016. ``APL in R.''
doi:\href{https://doi.org/10.13140/RG.2.1.2372.0724}{10.13140/RG.2.1.2372.0724}.


\end{document}
